% U9 WS5 -- nonstandard conditions, electrolysis, + 10 pt FRQ. Block 9.
\documentclass{apchem-worksheet}

\apcunit{9}
\apcunittitle{Thermodynamics and Electrochemistry}
\apcdoctitle{Worksheet 5 \textbullet\ Nonstandard Conditions, Electrolysis \& FRQ}
\apcsubtitle{CED 9.10--9.11 \textbullet\ argue from $Q$ against $K$; no Nernst equation required}

\begin{document}
\wsheader[$q = It$ \textbullet\ $\text{mol e}^- = q/F$ \textbullet\
$F = \SI{96485}{\coulomb\per\mole}$ \textbullet\
$E^\circ_{\text{cell}} = E^\circ_{\text{cathode}} -
E^\circ_{\text{anode}}$ \textbullet\ $\Delta G^\circ = -nFE^\circ$
\textbullet\ $E^\circ$ (V): \ce{Ag+}/\ce{Ag} $+0.80$ \textbullet\
\ce{Cu^2+}/\ce{Cu} $+0.34$ \textbullet\ \ce{Zn^2+}/\ce{Zn} $-0.76$
\textbullet\ molar masses: \ce{Ag} 107.87, \ce{Cu} 63.55, \ce{Al} 26.98.]

\problem[]{\ced{9.10} A cell's voltage as it runs.}
\begin{parts}
  \item As a galvanic cell operates, $Q$ \answerblank[1.6cm]{rises} toward
        $K$ and $E_{\text{cell}}$ \answerblank[1.6cm]{falls}.
  \item The cell is dead when \answerblank[3cm]{$Q = K$}, at which point
        $E_{\text{cell}} = $ \answerblank[1.4cm]{0}.
  \item Explain in one sentence what a cell potential physically measures.
        \answerlines[2]{How far the cell reaction currently sits from
        equilibrium --- the larger the displacement, the greater the
        driving force available to push electrons through the wire.}
\end{parts}

\problem[]{\ced{9.10} For \ce{Zn(s) + Cu^2+(aq) -> Zn^2+(aq) + Cu(s)},
predict the effect on $E_{\text{cell}}$ and justify each with $Q$. Do not
calculate.}
\begin{parts}
  \item $[\ce{Cu^2+}]$ is increased.
        \answerblank[6.6cm]{increases --- $Q$ falls}
  \item $[\ce{Zn^2+}]$ is increased.
        \answerblank[6.6cm]{decreases --- $Q$ rises}
  \item More solid \ce{Zn} is added.
        \answerblank[7.2cm]{no change --- solids omitted from $Q$}
  \item Write the reasoning behind (a) in full, as an FRQ answer would
        require. \answerlines[3]{\ce{Cu^2+} appears in the denominator of
        $Q$, so raising its concentration lowers $Q$. A smaller $Q$ places
        the system further from equilibrium, which increases the driving
        force, so $E_{\text{cell}}$ increases.}
\end{parts}

\problem[]{\ced{9.10} Concentration cells. Two \ce{Cu} electrodes sit in
\SI{0.10}{\Molar} and \SI{1.0}{\Molar} \ce{Cu^2+}.}
\begin{parts}
  \item $E^\circ_{\text{cell}}$ for this cell is
        \answerblank[1.4cm]{0}, because
        \answerblank[6cm]{both half-cells are identical}.
  \item Yet it produces a voltage. Which half-cell is the anode, and why?
        \answerlines[3]{The dilute (\SI{0.10}{\Molar}) half-cell. Copper
        dissolves there, raising its \ce{Cu^2+} concentration, while
        \ce{Cu^2+} plates out in the concentrated half-cell and lowers it.
        The cell runs in whichever direction moves the two concentrations
        toward each other.}
  \item When does it stop? \answerblank[6.9cm]{when the concentrations are
        equal}
\end{parts}

\problem[]{\ced{9.11} Faraday's law --- the chain. Put the four conversions
in order:}
\begin{parts}
  \item First: \answerblank[5.5cm]{current and time to charge}
  \item Then: \answerblank[5.7cm]{charge to moles of electrons}
  \item Then: \answerblank[8.1cm]{moles of electrons to moles of
        substance}
  \item Finally: \answerblank[5cm]{moles to grams}
  \item Which step needs the half-reaction, and why?
        \answerlines[2]{The third. The half-reaction gives the number of
        electrons per formula unit --- one for \ce{Ag+}, two for
        \ce{Cu^2+}, three for \ce{Al^3+} --- and that ratio cannot be
        guessed from the current.}
\end{parts}

\problem[]{\ced{9.11} A current of \SI{2.00}{\ampere} flows for
\SI{30.0}{\minute} through \ce{AgNO3}. Calculate the mass of silver
deposited. \workspace[2.8cm]
\keyonly{\begin{apcanswer}$t = 1800$~s; $q = (2.00)(1800) = 3600$~C;
mol e$^- = 3600/96485 = 0.0373$; \ce{Ag+ + e^- -> Ag} is one electron per
atom, so $m = (0.0373)(107.87) = \mathbf{4.03}$~g\end{apcanswer}}}

\problem[]{\ced{9.11} The same charge is passed through \ce{CuSO4} instead
--- \SI{1.50}{\ampere} for \SI{1.00}{\hour}.}
\begin{parts}
  \item Calculate the mass of copper deposited. \workspace[2.6cm]
        \keyonly{\begin{apcanswer}$q = (1.50)(3600) = 5400$~C;
        mol e$^- = 5400/96485 = 0.0560$; \ce{Cu^2+ + 2e^- -> Cu} needs two
        per atom, so mol Cu $= 0.0280$ and
        $m = (0.0280)(63.55) = \mathbf{1.78}$~g\end{apcanswer}}
  \item A student reports 3.56 g. Identify the single omitted step.
        \answerblank[6.4cm]{did not divide by $n = 2$}
\end{parts}

\problem[]{\ced{9.11} How long, in hours, to deposit \SI{5.00}{\gram} of
copper at \SI{2.00}{\ampere}? \workspace[2.8cm]
\keyonly{\begin{apcanswer}mol Cu $= 5.00/63.55 = 0.0787$;
mol e$^- = 2(0.0787) = 0.157$; $q = (0.157)(96485) = 15{,}183$~C;
$t = 15{,}183/2.00 = 7591$~s $= \mathbf{2.11}$~h\end{apcanswer}}}

\problem[]{\ced{9.11} Producing aluminium requires three electrons per
atom and \ce{Al^3+}/\ce{Al} sits at $-1.66$~V. Explain why aluminium
smelters are built beside power stations, and why recycling saves so much
energy.
\answerlines[3]{Reducing \ce{Al^3+} is strongly unfavorable, so a large
applied potential is needed, and three moles of electrons are required per
mole of metal. Both factors make the electrical energy per kilogram very
large. Recycling melts aluminium that has already been reduced, skipping
the electrolysis entirely.}}

\problem[]{\textbf{FRQ (10 points).} \ced{9.8} \ced{9.9} \ced{9.10}
A galvanic cell is constructed from a zinc electrode in
\SI{1.0}{\Molar} \ce{Zn(NO3)2} and a silver electrode in \SI{1.0}{\Molar}
\ce{AgNO3}, joined by a salt bridge.}
\begin{parts}
  \item Identify the anode and the cathode, and state the direction of
        electron flow through the external wire.
        \answerlines[2]{Zinc is the anode (oxidation, the more negative
        reduction potential); silver is the cathode (reduction). Electrons
        flow through the wire from the zinc electrode to the silver
        electrode.}
  \item Write the balanced overall cell reaction and give $n$.
        \answerlines[2]{\ce{Zn(s) + 2Ag+(aq) -> Zn^2+(aq) + 2Ag(s)};
        $n = 2$.}
  \item Calculate $E^\circ_{\text{cell}}$. \workspace[1.8cm]
        \keyonly{\begin{apcanswer}$E^\circ = 0.80 - (-0.76) =
        \mathbf{+1.56}$~V\end{apcanswer}}
  \item Calculate $\Delta G^\circ$ in kJ. \workspace[2.4cm]
        \keyonly{\begin{apcanswer}$\Delta G^\circ = -nFE^\circ =
        -(2)(96485)(1.56) = -301{,}033~\text{J} =
        \mathbf{-301}~\text{kJ}$\end{apcanswer}}
  \item The \ce{AgNO3} solution is diluted to \SI{0.10}{\Molar}. Predict
        the effect on $E_{\text{cell}}$ and justify without calculating.
        \answerlines[3]{\ce{Ag+} is a reactant, so lowering its
        concentration raises $Q$. A larger $Q$ places the system closer to
        equilibrium, reducing the driving force, so $E_{\text{cell}}$
        decreases.}
  \item State the value of $E_{\text{cell}}$ when the cell is exhausted,
        and the relationship between $Q$ and $K$ at that moment.
        \answerblank[7.4cm]{$E_{\text{cell}} = 0$ and $Q = K$}
\end{parts}

\begin{apcnote}[Rubric --- score yourself, 10 points]
\textbf{(a) 2 pts} 1 for both electrodes correctly assigned, 1 for electron
direction stated as anode $\to$ cathode \textbullet\ \textbf{(b) 2 pts}
1 balanced equation with \ce{Ag+} doubled, 1 for $n = 2$ --- an unbalanced
equation earns 0 for both \textbullet\ \textbf{(c) 1 pt} $+1.56$~V; using
$2(0.80)$ earns 0 \textbullet\ \textbf{(d) 2 pts} 1 correct substitution
into $-nFE^\circ$ with $n = 2$, 1 answer converted to kJ --- reporting
joules as kilojoules earns 0 for that point \textbullet\ \textbf{(e) 2 pts}
1 states that $Q$ increases because \ce{Ag+} is a reactant, 1 concludes
$E_{\text{cell}}$ decreases --- a bare ``decreases'' with no $Q$ reasoning
earns 0 of 2 \textbullet\ \textbf{(f) 1 pt} both $E = 0$ and $Q = K$
required.

\textbf{Total: } $2 + 2 + 1 + 2 + 2 + 1 = 10$.
\end{apcnote}

\begin{spiralstrip}{Unit 9 \textbullet\ blocks 1--8}
  \item $E^\circ_{\text{cell}}$ for \ce{Zn} and \ce{Ag+}:
        \answerblank[2.6cm]{$+1.56$ V}
  \item $\Delta G^\circ = $ \answerblank[3cm]{$-nFE^\circ$}
  \item Oxidation occurs at the: \answerblank[2cm]{anode}
  \item Crossover temperature $= $ \answerblank[3cm]{$\Delta H^\circ /
        \Delta S^\circ$}
  \item A catalyst leaves which two unchanged?
        \answerblank[4cm]{$\Delta G^\circ$ and $K$}
\end{spiralstrip}

\end{document}
